\section*{الفصل \ref{ch:g4:solar-system} --- المجموعة الشمسية}

\begin{solution}{exo:g4:solar-system:1}
\omterm{def:g4:solar-system:star}{النجم} يصنع \omterm{def:g1:light-and-shadows:source}{ضوءه} بنفسه؛ \omterm{def:g4:solar-system:planet}{والكوكب} لا يصنع ضوءًا ويسافر حول \omterm{def:g4:solar-system:star}{نجم}،
ولا يسطع إلا \omterm{def:g1:light-and-shadows:source}{بضوء} شمس مردود. فالشمس \omterm{def:g4:solar-system:star}{نجم}؛ والأرض \omterm{def:g4:solar-system:planet}{كوكب}.
\end{solution}

\begin{solution}{exo:g4:solar-system:2}
عطارد، والزهرة، والأرض، والمريخ، والمشتري، وزحل، وأورانوس،
ونبتون --- وكوكبنا هو الثالث، الأرض.
\end{solution}

\begin{solution}{exo:g4:solar-system:3}
الصخرية: عطارد والزهرة والأرض والمريخ. والعمالقة: المشتري وزحل
وأورانوس ونبتون. والأكبر: المشتري. وصاحب الحلقات: زحل.
\end{solution}

\begin{solution}{exo:g4:solar-system:4}
لأنه يدور حول الأرض لا حول الشمس مباشرة --- فهو \omterm{def:g2:moon-in-the-sky:moon}{قمر}، رفيق
\omterm{def:g4:solar-system:planet}{لكوكب}، لا \omterm{def:g4:solar-system:planet}{كوكب}.
\end{solution}

\begin{solution}{exo:g4:solar-system:5}
كلما بَعُد الطريق عن الشمس طالت الدورة: فعطارد القريب يسابق
سريعًا، ونبتون البعيد يستغرق أعمارًا.
\end{solution}

\begin{solution}{exo:g4:solar-system:6}
كرة القدم هي الشمس؛ وتقف الأرض (وهي حبّة فلفل في النموذج) على
نحو $30$ خطوة كبيرة. \omterm{def:g4:solar-system:family}{والمجموعة الشمسية} من حيث الحيّز فضاء خالٍ
كلها تقريبًا.
\end{solution}

\begin{solution}{exo:g4:solar-system:7}
من الشمس: فالزهرة، مثل \omterm{def:g2:moon-in-the-sky:moon}{القمر}، لا تفعل إلا أن تردّ \omterm{def:g1:light-and-shadows:source}{ضوء} الشمس ---
وغيومها البيضاء تردّه ردًّا سخيًّا حتى تفوق سطوعًا كل \omterm{def:g4:solar-system:star}{نجم}
حقيقي في سماء مسائنا.
\end{solution}

\begin{solution}{exo:g4:solar-system:8}
راقبه في مواجهة البريقات الثابتة بضعة أسابيع: \omterm{def:g4:solar-system:planet}{فالكوكب} ينزاح على
مهل --- يهيم --- بينها، أما النجوم الحقيقية فتحفظ مواضعها في
الأشكال.
\end{solution}

\begin{solution}{exo:g4:solar-system:9}
$150 \div 30 = 5$: فالمشتري أبعد عن الشمس من الأرض بنحو خمس
مرات.
\end{solution}

\begin{solution}{exo:g4:solar-system:10}
تكذب كل الرسوم تقريبًا في المسافات: فلو رُسمت بمقياس واحد
لتقلّصت الكواكب في صفحة تُظهر مدار نبتون حتى تصير نقاطًا لا
تُرى --- فيحشر الرسّامون المدارات حتى يُظهروا العوالِم أصلًا.
والحقيقة --- كما يبيّن النموذج الممشيّ --- أن الكواكب تسافر في
خواء واسع: حتى \omterm{def:g1:light-and-shadows:source}{الضوء} والمسبارات السريعة عليها أن تقطع مئات
``الخطوات'' من اللاشيء، وذلك يستغرق سنوات.
\end{solution}

\begin{solution}{exo:g4:solar-system:11}
الهائمات تنزاح أحيانًا إلى الوراء بين النجوم أسابيع --- ويثقل
تفسير ذلك وكل شيء يدور حول أرض متوّجة. أما والشمس في المركز
فالانزياح إلى الوراء تجاوزٌ بسيط: فالأرض، على طريق داخلي أسرع،
تتجاوز كوكبًا خارجيًّا أبطأ، فيبدو ذلك \omterm{def:g4:solar-system:planet}{الكوكب} منزلقًا إلى
الوراء. والثمن: التخلّي عن المركز --- وقبول أن بيتنا هو \omterm{def:g4:solar-system:planet}{الكوكب}
رقم ثلاثة.
\end{solution}
